\documentclass[12pt]{report}
\usepackage[top=2.3cm,bottom=2.5cm,left=2.5cm,right=2.5cm]{geometry}
\usepackage[T1]{fontenc}
\usepackage{lmodern}
\usepackage{microtype}
\usepackage{graphicx}
\usepackage{booktabs}
\usepackage{enumitem}
\setlist{itemsep=0.15em,topsep=0.35em,parsep=0pt}
\usepackage{array}
\usepackage{float}
\usepackage{longtable}
\usepackage{tabularx}
\usepackage{amsmath}
\usepackage{amssymb}
\usepackage{parskip}
\usepackage[hidelinks]{hyperref}

\begin{document}
\pagestyle{plain}
\setcounter{chapter}{0}
\chapter{Introduction to Artificial Intelligence}

\section{Introduction to Artificial Intelligence}

Artificial Intelligence (AI) represents one of the most transformative disciplines in modern computer science. It focuses on constructing computational systems capable of performing tasks that traditionally require human intelligence, such as logical reasoning, decision-making, visual perception, and natural language understanding. This chapter provides a comprehensive foundation for CSC 309 by exploring the historical evolution, core principles, operational methodologies, and major theoretical debates that define artificial intelligence.

Understanding the conceptual and historical origins of AI is essential for contextualizing modern computational approaches. By examining how symbolic reasoning gave way to probabilistic modeling and data-driven learning, students gain insight into why modern AI architectures are designed in specific ways. Furthermore, evaluating theoretical frameworks such as the Turing test and John Searle's Chinese Room thought experiment equips students with the critical analysis skills needed to evaluate system capabilities, performance boundaries, and ethical implications.

This introductory chapter serves as the foundation for the entire CSC 309 curriculum. The concepts, techniques, and taxonomy established here directly connect to subsequent chapters on intelligent agents, state-space search algorithms, formal knowledge representation, Natural Language Processing, expert systems, and robotics.

\subsection{Definition and Primary Goals of Artificial Intelligence}

\textbf{Artificial Intelligence} is defined as a subfield of computer science dedicated to creating hardware and software systems capable of performing tasks that typically demand human cognitive abilities, including learning, reasoning, problem-solving, perception, and language comprehension.

The field of AI is driven by two primary goals:
\begin{enumerate}
    \item \textbf{Engineering Goal}: To solve complex, real-world problems by designing, building, and deploying reliable, autonomous, and intelligent algorithms.
    \item \textbf{Scientific Goal}: To model, formalize, and understand the mechanisms underlying cognitive processes and human intelligence through computational simulation.
\end{enumerate}

To achieve these goals, artificial intelligence encompasses several specific sub-goals:
\begin{itemize}
    \item Logical reasoning and automated decision-making under uncertain or partial information.
    \item Autonomous knowledge acquisition and continuous performance improvement through learning algorithms.
    \item Perceptual capabilities involving visual, auditory, and sensor-based environmental processing.
    \item Natural language interaction and semantic comprehension for seamless communication.
\end{itemize}

\subsection{Historical Evolution of Artificial Intelligence}

The historical trajectory of artificial intelligence spans several distinct eras, characterized by shifting paradigms, periods of high optimism, and subsequent periods of reduced funding known as AI winters.

\begin{verbatim}
[1940s-1950]  Foundational Roots (Turing's Imitation Game)
      |
[1956]        Dartmouth Workshop (Birth of Artificial Intelligence)
      |
[1956-1970s]  Early Enthusiasm & Symbolic AI (GPS, ELIZA)
      |
[1974-1980]   First AI Winter (Lighthill Report, combinatorial explosion)
      |
[1980-1987]   Expert Systems Boom (Knowledge-based rule engines)
      |
[1987-1993]   Second AI Winter (Hardware collapse, maintenance costs)
      |
[1990s-Pres]  Modern Machine Learning & Deep Learning (Data-driven AI)
\end{verbatim}

\begin{enumerate}
    \item \textbf{Foundational Roots (1940s--1950)}: Alan Turing laid the theoretical groundwork for machine intelligence. In his seminal 1950 paper, "Computing Machinery and Intelligence," Turing posed the fundamental question, "Can machines think?" and introduced the operational framework known as the Imitation Game.
    \item \textbf{The Dartmouth Workshop (1956)}: The formal birth of AI as an academic discipline occurred at the Dartmouth Summer Research Project on Artificial Intelligence organized by John McCarthy, Marvin Minsky, Nathaniel Rochester, and Claude Shannon. John McCarthy coined the term \textbf{Artificial Intelligence} during this historic gathering.
    \item \textbf{Early Enthusiasm and Symbolic AI (1956--1970s)}: Early AI research focused on symbolic logic, formal theorem proving, and micro-world problem solving. Systems such as the General Problem Solver (GPS) and ELIZA demonstrated early potential, leading to high optimism.
    \item \textbf{The First AI Winter (1974--1980)}: Early systems struggled with real-world complexity due to computational hardware constraints and combinatorial explosion. Critical reports, such as the Lighthill Report in the United Kingdom, led to widespread reductions in research funding.
    \item \textbf{The Expert Systems Boom (1980--1987)}: Industrial interest re-emerged with the adoption of knowledge-based systems. \textbf{Expert systems} used structured rules to emulate human domain expertise in specialized areas such as medical diagnostics and geological exploration.
    \item \textbf{The Second AI Winter (1987--1993)}: Specialized hardware markets collapsed, and rule-based expert systems proved expensive to maintain and difficult to scale to unmodeled scenarios.
    \item \textbf{The Modern Era of Machine Learning and Deep Learning (1990s--Present)}: Exponential increases in processing power, high-throughput memory architectures, and large datasets propelled a shift from hand-coded logical rules to statistical data-driven learning. Deep neural network architectures enabled breakthroughs across computer vision, speech recognition, and Natural Language Processing.
\end{enumerate}


\section{Fundamental AI Techniques and Taxonomy}

\subsection{Core AI Techniques}

AI applications rely on four fundamental technical paradigms:

\begin{enumerate}
    \item \textbf{Knowledge Representation and Logic}: Structuring domain facts, rules, and relationships into formal computational structures that permit automated inference and deduction.
    \item \textbf{Search and Optimization}: Navigating vast state spaces systematically to identify optimal or satisfying action sequences using search algorithms and explicit evaluation metrics.
    \item \textbf{Machine Learning and Statistical Pattern Recognition}: Applying mathematical models that automatically infer patterns, underlying distributions, and decision boundaries from training data without explicit manual programming.
    \item \textbf{Neural Networks and Deep Learning}: Processing complex high-dimensional data through layered computational graphs that extract hierarchical representations directly from raw inputs.
\end{enumerate}

\subsection{Classification and Types of Artificial Intelligence}

Artificial intelligence systems are classified according to their capability levels and functional architectures.

\subsubsection{Classification by Capability Level}

\begin{verbatim}
+-----------------------------------------------------------------+
|               Artificial Superintelligence (ASI)                |
|       (Surpasses human intellect across all cognitive domains)  |
+-----------------------------------------------------------------+
                                ^
                                |
+-----------------------------------------------------------------+
|                Artificial General Intelligence (AGI)             |
|       (Human-equivalent cross-domain reasoning & learning)      |
+-----------------------------------------------------------------+
                                ^
                                |
+-----------------------------------------------------------------+
|                Artificial Narrow Intelligence (ANI)             |
|       (Specialized single-domain execution; all current AI)     |
+-----------------------------------------------------------------+
\end{verbatim}

\begin{itemize}
    \item \textbf{Artificial Narrow Intelligence (ANI)}: Also referred to as Weak AI, ANI describes systems built and optimized to perform a specific, isolated task (e.g., chess software, spam filters, or image recognition models). ANI systems cannot transfer their functional expertise to unassigned tasks. All operational AI systems today belong to this category.
    \item \textbf{Artificial General Intelligence (AGI)}: Also known as Strong AI, AGI refers to theoretical systems possessing human-level cognitive flexibility. An AGI system could perform cross-domain reasoning, transfer learning across disparate subjects, and adapt dynamically to unfamiliar environments.
    \item \textbf{Artificial Superintelligence (ASI)}: A theoretical extension of AGI where machine cognitive abilities surpass human intellect across all computational, creative, scientific, and social domains.
\end{itemize}

\subsubsection{Classification by Functional Architecture}

\begin{itemize}
    \item \textbf{Reactive Machines}: Basic systems that calculate actions purely from current environmental inputs without maintaining internal memory or past state history (e.g., IBM's Deep Blue).
    \item \textbf{Limited Memory}: Systems capable of retaining dynamic short-term memory to inform real-time decision-making (e.g., trajectory planning in autonomous vehicle navigation).
    \item \textbf{Theory of Mind}: Theoretical frameworks capable of representing and reasoning about the psychological states, intentions, beliefs, and emotions of human agents.
    \item \textbf{Self-Aware AI}: Theoretical systems possessing explicit self-consciousness, self-identity, and self-referential mental monitoring.
\end{itemize}

\begin{table}[htbp]
\centering
\begin{tabularx}{\textwidth}{@{} >{\raggedright\arraybackslash}p{3.5cm} X X @{}}
\toprule
\textbf{Category} & \textbf{Key Characteristics} & \textbf{Current Operational Status} \\
\midrule
\textbf{Artificial Narrow Intelligence (ANI)} & Programmed for specialized, single-domain execution; incapable of cross-domain transfer. & Deployed universally in production (e.g., search engines, diagnostics). \\
\textbf{Artificial General Intelligence (AGI)} & Human-equivalent cognitive versatility; adapts autonomously across open domains. & Theoretical research objective; no physical instances exist. \\
\textbf{Artificial Superintelligence (ASI)} & Cognitive capacity exceeding human capability across all problem domains. & Hypothetical concept; focus of safety and ethics research. \\
\bottomrule
\end{tabularx}
\end{table}


\section{Evaluating Machine Intelligence}

\subsection{The Turing Test}

The \textbf{Turing test}, proposed by Alan Turing in 1950, provides an operational methodology for evaluating whether a system can display behavior indistinguishable from human intelligence.

In the standard operational setup, a human interrogator ($C$) communicates with two hidden entities through a text-based terminal interface: a human ($B$) and a machine ($A$). The interrogator submits text queries and evaluates the text responses. The machine's goal is to deceive the interrogator into believing it is the human. If the interrogator cannot reliably distinguish the machine from the human after a specified evaluation duration, the machine passes the test.

\begin{verbatim}
+-------------------------------------------------+
|                  Interrogator (C)              |
+-------------------------------------------------+
                        |
            +-----------+-----------+
            | Text-Based Interface  |
            +-----------+-----------+
                        |
       +----------------+----------------+
       |                                 |
       v                                 v
+--------------+                  +--------------+
| Machine (A)  |                  |  Human (B)   |
+--------------+                  +--------------+
\end{verbatim}

Despite its historical significance, the Turing test faces critical theoretical objections:
\begin{itemize}
    \item It measures external behavioral imitation rather than underlying mental cognition or comprehension.
    \item Simple trickery, pre-programmed evasion techniques, or superficial surface matching can deceive human interrogators.
    \item It is anthropocentric, evaluating intelligence strictly through the lens of human linguistic behavior.
\end{itemize}

\subsection{The Chinese Room Thought Experiment}

Formulated by philosopher John Searle in 1980, the \textbf{Chinese Room thought experiment} challenges the claim that functional symbol manipulation constitutes true cognitive understanding or intentionality.

Searle asks us to imagine a monolingual English speaker who knows no Chinese, locked inside a sealed room. The room contains a structured rulebook written in English. Outside observers slide slips of paper containing Chinese symbols into the room through a slot. The occupant uses the English rulebook to look up incoming Chinese character sequences and follow exact instructions for matching them to corresponding output Chinese characters. The occupant then slides the output slips back out through a second slot.

\begin{verbatim}
Input: Chinese Symbols
        |
        v
+-------------------------------------------------------+
|                    The Chinese Room                   |
|                                                       |
|   +-------------------+       +-------------------+   |
|   | English Monolingual|       |  Rulebook/Program |   |
|   |  Symbol Operator  | <---> | (Syntactic Matching|   |
|   +-------------------+       +-------------------+   |
+-------------------------------------------------------+
        |
        v
Output: Chinese Symbols
\end{verbatim}

To an outside Chinese speaker, the responses emerging from the room are coherent and intelligent. However, the person inside the room processes the symbols purely based on syntactic rules without understanding any semantic meaning of the Chinese characters.

Searle uses this experiment to make a key distinction:
\begin{itemize}
    \item \textbf{Syntax}: Formal rules governing symbol manipulation and computational structure.
    \item \textbf{Semantics}: The actual meaning, intentionality, and underlying understanding associated with symbols.
\end{itemize}

Searle concludes that digital computers operate entirely syntactically. Consequently, running a program, regardless of its operational complexity, cannot produce genuine semantic understanding, mental states, or intentionality.

\begin{table}[htbp]
\centering
\begin{tabularx}{\textwidth}{@{} >{\raggedright\arraybackslash}p{3cm} X X @{}}
\toprule
\textbf{Dimension} & \textbf{Turing Test} & \textbf{Chinese Room Thought Experiment} \\
\midrule
\textbf{Primary Proponent} & Alan Turing (1950) & John Searle (1980) \\
\textbf{Core Criterion} & External operational output and conversational capability. & Internal semantic state and conscious intentionality. \\
\textbf{Main Conclusion} & Behavior indistinguishable from human intelligence implies cognitive capability. & Correct symbol manipulation can occur entirely without semantic understanding. \\
\textbf{Philosophical View} & Functionalism / Behaviorism. & Biological Naturalism / Anti-functionalism. \\
\bottomrule
\end{tabularx}
\end{table}


\section{Branches, Applications, and Operational Trade-Offs}

\subsection{Key Branches and Real-World Applications}

Artificial intelligence spans multiple primary subfields:

\begin{itemize}
    \item \textbf{Natural Language Processing (NLP)}: Algorithms designed for processing, understanding, translating, and generating natural human language text and speech.
    \item \textbf{Computer Vision}: Techniques for capturing, transforming, and interpreting visual information from digital image streams or camera arrays.
    \item \textbf{Robotics}: Integration of perceptual sensors, decision-making controllers, and physical actuators to carry out tasks in dynamic real-world environments.
    \item \textbf{Expert Systems}: Knowledge-based software engines that use rule structures to simulate domain-specific human decision-making.
    \item \textbf{Machine Learning}: Statistical modeling techniques that enable systems to improve task performance automatically through data experience.
\end{itemize}

AI technologies are deployed across major industrial domains:
\begin{itemize}
    \item \textbf{Healthcare}: Automated diagnostic image interpretation, personalized genomic analysis, and predictive clinical risk modeling.
    \item \textbf{Finance}: High-frequency algorithmic stock trading, real-time banking transaction fraud detection, and credit assessment models.
    \item \textbf{Transportation}: Autonomous vehicle obstacle avoidance, dynamic urban traffic lights regulation, and supply chain path optimization.
    \item \textbf{E-Commerce}: Dynamic user recommendation engines, inventory management, and automated customer support interfaces.
\end{itemize}

\subsection{Advantages and Disadvantages of Artificial Intelligence}

Deploying artificial intelligence systems involves significant operational advantages alongside structural risks.

\begin{table}[htbp]
\centering
\begin{tabularx}{\textwidth}{@{} X X @{}}
\toprule
\textbf{Advantages} & \textbf{Disadvantages and Risks} \\
\midrule
Continuous 24/7 operational availability without fatigue or performance degradation. & High initial capital expenditure, ongoing infrastructure costs, and energy consumption. \\
High-speed data evaluation and statistical accuracy across massive datasets. & System susceptibility to algorithmic bias inherited from skewed training data. \\
Automation of hazardous, repetitive, or error-prone tasks. & Structural workforce dislocation and job displacement across various sectors. \\
Consistent execution of complex rules without human variance. & Complete absence of moral judgment, human empathy, and genuine common-sense context. \\
\bottomrule
\end{tabularx}
\end{table}


\section{Conclusion}

This chapter established the foundational history, core techniques, functional taxonomies, and theoretical questions surrounding artificial intelligence. While empirical metrics such as the Turing test and arguments like the Chinese Room thought experiment explore the boundary between behavioral simulation and genuine cognitive understanding, practical AI systems rely on targeted techniques, such as search algorithms, knowledge representation, and statistical machine learning, to solve real-world problems. 

In Chapter 2, we formalize how these individual AI capabilities are unified into structured systems known as \textbf{intelligent agents}, analyzing how agents interact with diverse operational environments to maximize performance metrics.


\section*{Tutorial Questions}

\begin{enumerate}
    \item Distinguish between the engineering goal and the scientific goal of artificial intelligence. Provide a practical application scenario illustrating each goal.
    \item Outline the primary factors that caused the first AI winter during the 1970s. How did the transition from symbolic AI to data-driven machine learning address these systemic limitations?
    \item Compare Artificial Narrow Intelligence (ANI), Artificial General Intelligence (AGI), and Artificial Superintelligence (ASI). Why are all current operational AI systems classified as ANI?
    \item Describe the operational configuration of the Turing test. Discuss two major conceptual criticisms of using the Turing test as a definitive metric for genuine intelligence.
    \item Explain John Searle's Chinese Room thought experiment. How does Searle use this experiment to differentiate between syntax and semantics, and what are its implications for Strong AI?
    \item Compare reactive machines with limited memory AI systems, giving one concrete example of each system type.
    \item An automated diagnostic system is trained on historical hospital records to flag high-risk patient admissions. Identify two operational advantages and two technical or ethical risks associated with deploying this AI system.
    \item Explain why rule-based expert systems struggled to scale when applied to complex real-world environments, and describe how statistical machine learning techniques overcome these scalability barriers.
\end{enumerate}

\end{document}
